\par\SourceLeaf{528}The rule can therefore also be expressed as follows; (always assuming $p\equiv5\ (\mathrm{mod}.8)$)
\[
\begin{aligned}
\text{We have }b&\equiv a+1\ (\mathrm{mod}.8),&\quad&\text{if }m\text{ is even}\\
b&\equiv a+5&&\text{if }m\text{ is odd.}
\end{aligned}
\]
I do not yet venture a conjecture as to whether a still simpler criterion is possible for distinguishing the case of even $m$ from the odd case \textit{in advance}, that is, without knowing the value of $m$ itself, since, as I remarked above, this connection is still entirely new.
\newpage
\par For the case $p\equiv1\ (\mathrm{mod}.8)$ the above congruence $b\equiv2k+a-1\ (\mathrm{mod}.8)$ remains correct, but it no longer determines the sign of $b$, since both the positive and the negative values of $b$ satisfy it. Here $k$ is always even, $=2m$ (with $m$ understood just as above), or $=2m+2$ if $m$ denotes the number of decompositions of $p$ into 3 positive \textit{unequal} squares, and $b\equiv0\ (\mathrm{mod}.4)$, or $b\equiv-b\ (\mathrm{mod}.8)$. I conjecture that the case $p\equiv1\ (\mathrm{mod}.8)$ or $b\equiv0\ (\mathrm{mod}.4)$ calls for deeper investigation, and perhaps again
\[
\begin{aligned}
b&\equiv4\ (\mathrm{mod}.8)&&\text{easier than }b\equiv0\ (\mathrm{mod}.8)\\
b&\equiv8\ (\mathrm{mod}.16)&&\text{easier than }b\equiv0\ (\mathrm{mod}.16)
\end{aligned}
\]
etc.
\par\Needspace{5\baselineskip}With the highest esteem I remain
\par\nobreak\noindent\hfill Yours most devotedly in friendship
\par\nobreak\noindent Göttingen, 30 May 1828.\hfill C. F. \textsc{Gauss}.
\BookRule{.25\textwidth}
